\section{Every odd modulus}
\label{sec:odd-modulus-abel}
Every odd modulus \(N\ge3\) has pure Abel value strictly above \(3/4\) or
strictly below \(5/7\). The canonical proof is
\path{mathematics/fixed-map/odd_modulus_abel.md}.
The integer comparisons are certified by
\path{formalization/Bgs/OddModulusRational.lean} and recomputed by the
Python and C++ verifiers named there. No external refereeing is asserted.
BGS-TRACE-005 and BGS-BGS-001 remain open.

\begin{theorem}[Odd moduli above \(3/4\) or below \(5/7\), BGS-TRACE-124]
\label{thm:odd-modulus-abel}
Let \(N\ge3\) be odd. Exactly one of the following holds:
\(L_N(M_N;1)>3/4\), so the distance from this value to
\((1-e^{-2})/4\) is greater than \(1/2\), or
\(L_N(M_N;1)<5/7\), so that distance is less than \(1/2\).
The modulus \(315=3^2\cdot5\cdot7\) is in the first case, because its
radical \(105\) is short.

For the fixed kick \(\tau=10^{-80}\), the gap \(L_N(U_N;1)>3/4\) holds
when \(8\,t(N)\le3N\) and \(N\le10^{39}\), in particular at \(N=315\).
It holds when \(N\le10^{26}\) and \(N\) is a ratio-preserving tower in a
family whose pure margin exceeds \(13/420\): a split full-order tower over
a multiple of \(7\), the tower \(3\cdot7^a\) with \(a\ge2\), an inert
full-order tower whose exponent of \(7\) is at least \(2\), an inert
tower \(7q^b\) with exponent of \(7\) equal to \(1\) and
\(13\le q\le39\), a conditional non-short product of \(3\) with a short
prime of odd order, and \(3\cdot5^b\cdot13\) with \(b\ge2\). It holds when
\(N=15r\) with \(r\) an odd prime other than \(5\) and \(N\le2\cdot10^{26}\).

The kick is not claimed for a ratio-preserving tower over a ceiling
semiprime. At \(3\cdot23\) the gap from the tower majorant to \(5/7\) is
\(575/22594\), which is smaller than \(13/420\). The kick is not claimed
for an inert tower \(7q^b\) with the exponent of \(7\) equal to \(1\) and
\(q\ge40\): the comparison which proves \(L>3/4\) on that tail lies below
\(82/105\). The kick is not claimed for \(N=15r>2\cdot10^{26}\), nor at
the \(79\)-digit dimension \(u_{207}\). No density is proved. The signed
off-diagonal sum, the limit, and every \(k>2\) stay open.
\end{theorem}
